\documentclass[../../main.tex]{subfiles}
\begin{document}

\begin{center}
  {\fontsize{48pt}{58pt}\selectfont \textbf{1988}\normalsize\textbf{年度}}
\end{center}

\begin{table}[htb]
  \centering
  \caption{出題テーマと難易度}
  \begin{tabular}{|c|p{9cm}|c|}
    \hline
    \textbf{問題番号} & \textbf{テーマ} & \textbf{難易度} \\
    \hline
    第1問 & 一次変換と象限の関係 & \\
    \hline
    第2問 & 正四面体の正射影の面積 & \\
    \hline
    第3問 & 三次曲線の平行移動と共有点 & \\
    \hline
    第4問 & 傾きが変わる折れ線の極限 & \\
    \hline
    第5問 & 放物面の影と円弧の長さ & \\
    \hline
    第6問 & 四面体の体積の最大値 & \\
    \hline
  \end{tabular}
\end{table}

\begin{table}[htb]
  \centering
  \caption{解答}
  \renewcommand{\arraystretch}{1.4}
  \begin{tabular}{|c|c|p{9cm}|}
    \hline
    \textbf{問題番号} & \textbf{小問} & \textbf{答え} \\
    \hline
    第1問 & -- & \\
    \hline
    第2問 & -- & 最小値は$\dfrac{\sqrt{2}}{4}$，最大値は$\dfrac{1}{2}$ \\
    \hline
    第3問 & -- & 点$P$の領域は解答の図を参照 \\
    \hline
    第4問 & -- & $0<s\le1$，または$s>1$かつ$b<\dfrac{s}{a(s-1)}$ \\
    \hline
    第5問 & -- & 影の長さは$4\pi$ \\
    \hline
    第6問 & -- & 最大値は$9\sqrt{3}$ \\
    \hline
  \end{tabular}
\end{table}

\paragraph{試験形態}
6問．試験時間は要確認．

\end{document}
